\documentclass[12pt,a4paper]{article}
\newcommand{\CadernoDisciplina}{Análise Exploratória de Dados}
\newcommand{\CadernoCorHex}{17365D}
% Preâmbulo compartilhado dos cadernos de exercícios UNIFACOL.
% Definir \CadernoDisciplina e, opcionalmente, \CadernoCorHex antes deste input.
\providecommand{\CadernoDisciplina}{Disciplina}
\providecommand{\CadernoCorHex}{17365D}
\usepackage[a4paper,margin=2.05cm,headheight=15pt]{geometry}
\usepackage[T1]{fontenc}
\usepackage[utf8]{inputenc}
\usepackage[brazil]{babel}
\usepackage{lmodern,microtype,amsmath,amssymb}
\usepackage{enumitem,booktabs,tabularx,array,longtable,adjustbox,multirow}
\usepackage{xcolor,hyperref,graphicx,fancyhdr}
\usepackage[most]{tcolorbox}
\usepackage{tikz}
\usetikzlibrary{arrows.meta,positioning,fit,calc,patterns,shapes.geometric}
\definecolor{disciplinecolor}{HTML}{\CadernoCorHex}
\definecolor{stimulusgrey}{HTML}{EEF2F5}
\definecolor{answerorange}{HTML}{FFF0DE}
\definecolor{answerframe}{HTML}{B85C00}
\definecolor{supportgrey}{HTML}{F7F7F7}
\hypersetup{colorlinks=true,linkcolor=disciplinecolor,urlcolor=disciplinecolor}
\newcolumntype{Y}{>{\raggedright\arraybackslash}X}
\newcolumntype{C}{>{\centering\arraybackslash}X}
\setlength{\parindent}{0pt}
\setlength{\parskip}{.42em}
\setlist[enumerate]{leftmargin=*,itemsep=.28em,topsep=.35em}
\pagestyle{fancy}
\fancyhf{}
\lhead{UNIFACOL --- \CadernoDisciplina}
\rhead{Caderno ENADE --- 2026.2}
\cfoot{\thepage}
\newcommand{\ItemHeader}[4]{%
  \par\medskip\noindent
  \colorbox{disciplinecolor}{\parbox{\dimexpr\textwidth-2\fboxsep\relax}{\color{white}\bfseries Questão #1\hfill #2}}%
  \par\smallskip\noindent\textbf{Competência:} #3\par
  \noindent\textbf{Suporte:} #4\par\smallskip
}
\newcommand{\TextoBase}{\noindent\colorbox{stimulusgrey}{\parbox{\dimexpr\textwidth-2\fboxsep\relax}{\textbf{TEXTO-BASE}}}\par\smallskip}
\newcommand{\Comando}{\par\smallskip\noindent\textbf{COMANDO.} }
\newcommand{\FonteDidatica}{\par\smallskip\noindent{\footnotesize\textit{Fonte: situação e dados elaborados para fins didáticos.}}}
\newcommand{\FonteExterna}[1]{\par\smallskip\noindent{\footnotesize\textit{Fonte: #1}}}
\newcommand{\EspacoDiscursiva}{\par\noindent\rule{\textwidth}{.2pt}\vspace{1.15cm}\par\noindent\rule{\textwidth}{.2pt}}
\newtcolorbox{AnswerBox}{enhanced,breakable,colback=answerorange,colframe=answerframe,boxrule=.7pt,arc=1mm,title=Resposta explicada,fonttitle=\bfseries,coltitle=white,before skip=7pt,after skip=7pt}
\newtcolorbox{DocumentBox}[1][Documento ou artefato]{enhanced,breakable,colback=supportgrey,colframe=black!55,boxrule=.6pt,arc=.7mm,title={#1},fonttitle=\bfseries,coltitle=white,before skip=7pt,after skip=7pt}
\newtcolorbox{MemoBox}{enhanced,breakable,colback=white,colframe=disciplinecolor,boxrule=.7pt,arc=.7mm,title=Memorando,fonttitle=\bfseries,coltitle=white,before skip=7pt,after skip=7pt}
\newcommand{\LegendaDidatica}[1]{\par\smallskip{\footnotesize\textit{#1. Fonte: elaborado para fins didáticos.}}}

\setcounter{secnumdepth}{0}
\begin{document}
\begin{center}
{\Large\bfseries UNIFACOL}\\[2mm]
{\large Simulado ENADE}\\[1mm]
{\Large\bfseries VERSÃO B}
\end{center}
\begin{DocumentBox}[Identificação]
\begin{tabularx}{\textwidth}{@{}p{.48\textwidth}p{.48\textwidth}@{}}
Estudante: \hrulefill & Turma: \hrulefill\\[3mm]
Professor: Cayo Oliveira & Data: \rule{2.7cm}{.2pt}/2026
\end{tabularx}
\end{DocumentBox}
\textbf{Instruções.} Esta avaliação contém seis questões objetivas e duas discursivas. Nas objetivas, assinale uma única alternativa. Nas discursivas, mostre cálculos intermediários e justifique decisões com as evidências do texto-base. Respostas sem desenvolvimento podem não receber pontuação integral.
\bigskip
% origem: AED-C05-O09
\ItemHeader{1}{Objetiva}{Aplicar minimização por consumo}{SPEC de gestão}
\TextoBase A fonte acadêmica contém CPF, nome, curso, frequência e nota por matrícula. A SPEC executiva necessita apenas de taxas agregadas por curso e período; casos individuais são usados em acompanhamento pedagógico restrito. A publicação deve atender finalidade e acesso sem replicar identificadores ``para eventual uso futuro''.
\FonteDidatica
\Comando A publicação coerente é
\begin{enumerate}[label=\Alph*)]
  \item publicar agregados necessários na SPEC executiva e manter detalhe identificado apenas onde finalidade e acesso o justificam.
  \item aplicar hash no CPF e publicar todos os demais detalhes, pois hash elimina qualquer risco.
  \item incluir CPF para permitir futuras perguntas.
  \item remover curso e período, porque agregação impede decisão.
  \item copiar a SOR inteira e ocultar colunas na cor do painel.
\end{enumerate}

% origem: AED-C12-D15
\ItemHeader{2}{Discursiva}{Especificar teste de atualização}{Troca do modelo local}
\TextoBase
A equipe quer substituir Llama Q4 por uma versão nova e alterar simultaneamente prompt, parser e cache. O sistema atende perguntas numéricas, comparações e narrativas.
A atualização só poderá substituir a versão atual se passar casos funcionais, adversariais, de segurança e de regressão com configuração registrada.
\FonteDidatica
\Comando Proponha protocolo com baseline congelado, conjunto estratificado, versões, ablações, testes sintáticos e semânticos, cálculos de referência, latência, memória, custo, segurança, execução paralela, canário e rollback.
\EspacoDiscursiva

% origem: AED-C09-O05
\ItemHeader{3}{Objetiva}{contexto}{Situação profissional e suporte quantitativo}

\TextoBase

O histórico recuperado contém uma regra revogada e um dado atual, sem vigência visível. Se ambos forem enviados, o modelo pode produzir resposta fluente baseada na política antiga. O contexto precisa privilegiar fonte vigente e preservar proveniência.

\FonteDidatica

\Comando Selecione a estratégia de contexto que remove regra revogada sem perder vigência e proveniência.

\begin{enumerate}[label=\Alph*)]
  \item usar user para sobrescrever segurança
  \item enviar tudo porque mais contexto sempre melhora
  \item deixar o assistant decidir autoridade
  \item repetir regra antiga
  \item selecionar/versionar contexto, marcar proveniência e excluir ou rotular conteúdo revogado
\end{enumerate}

% origem: AED-C11-O07
\ItemHeader{4}{Objetiva}{Calcular capacidade e custo por resposta útil}{Teste de carga}
\TextoBase
Em 1.000 perguntas, uma configuração respondeu 900, das quais 810 foram corretas e citadas. Custou R\$180. Outra respondeu 800, das quais 760 foram corretas e citadas, ao custo de R\$190. O requisito é pelo menos 75\% de respostas úteis sobre o total e deseja-se menor custo por resposta útil.
A decisão de escala depende simultaneamente de cobertura útil e custo unitário, medidos sobre o mesmo conjunto de mil perguntas.
\FonteDidatica
\Comando A análise correta é
\begin{enumerate}[label=\Alph*)]
  \item somente a segunda passa, porque possui maior precisão entre as respostas geradas.
  \item a primeira alcança 81\% e R\$0,22 por útil; a segunda 76\% e R\$0,25; ambas passam, e a primeira tem menor custo unitário.
  \item a primeira alcança 90\% e R\$0,20; a segunda 80\% e R\$0,24; basta dividir custo por respostas geradas.
  \item somente a primeira passa, porque 800 respostas são menos que o limiar de 75\%.
  \item a primeira alcança 81\% e R\$0,18; a segunda 76\% e R\$0,19; custo deve ser dividido por perguntas totais.
\end{enumerate}

% origem: AED-C01-D11
\ItemHeader{5}{Discursiva}{Construir cadeia analítica verificável}{Caso de evasão}
\TextoBase Em solicitação de produto, uma faculdade propõe usar frequência, notas, situação financeira e contatos para ``prever evasão'' e enviar mensagens automáticas. O documento não define população, momento, resultado nem revisão humana, impedindo transformar a demanda em análise e intervenção avaliáveis.
\FonteDidatica
\Comando Em até 15 linhas, reformule a demanda como pergunta operacional; desenhe a cadeia pergunta--dado--análise--visual--decisão; indique dois riscos de qualidade ou ética e duas evidências necessárias antes de automatizar contato.
\EspacoDiscursiva

% origem: AED-C12-O03
\ItemHeader{6}{Objetiva}{Impedir execução arbitrária}{Pergunta analítica em linguagem natural}
\TextoBase
Para responder perguntas sobre um DataFrame, a equipe considera permitir que o Llama gere Python e execute qualquer código no processo do Streamlit, que também possui credenciais do banco e acesso à rede.
A pergunta pode conter texto malicioso na planilha; validação de intenção não substitui lista permitida de operações e isolamento da execução.
\FonteDidatica
\Comando A alternativa inicial mais segura e verificável é
\begin{enumerate}[label=\Alph*)]
  \item executar o código gerado e revisar apenas se ocorrer erro.
  \item oferecer ferramentas analíticas permitidas com parâmetros validados; se código livre for indispensável, isolá-lo em sandbox sem segredo, rede ou escrita ampla.
  \item ocultar as credenciais no prompt, mantendo-as disponíveis ao processo executor.
  \item usar temperatura zero como fronteira de segurança para a execução.
  \item pedir ao modelo uma promessa textual de que não acessará arquivos.
\end{enumerate}

% origem: AED-C01-O07
\ItemHeader{7}{Objetiva}{Distinguir dado, informação, insight e decisão}{Incidente de estoque}
\TextoBase Em uma ata de revisão, registros de venda são convertidos em 320 unidades por semana; a análise localiza rupturas sobretudo às sextas após promoções e recomenda antecipar reposição. A ata exige distinguir, nessa ordem, dado, informação, conhecimento e decisão monitorada.
\FonteDidatica
\Comando A classificação correta da cadeia é
\begin{enumerate}[label=\Alph*)]
  \item registros = dados; 320 = informação resumida; padrão contextualizado = insight; reposição monitorada = decisão.
  \item todos os elementos são dados, pois foram produzidos por um sistema.
  \item somente a reposição é informação, pois as etapas anteriores não geram ação.
  \item registros = decisão; 320 = dado bruto; padrão = ferramenta; reposição = informação.
  \item registros = insight; 320 = decisão; padrão = dado; reposição = visualização.
\end{enumerate}

% origem: AED-C07-O01
\ItemHeader{8}{Objetiva}{grão}{Situação profissional e suporte quantitativo}

\TextoBase

A tabela deveria ter uma linha por estudante-semestre, mas transferências criam duas linhas para a mesma pessoa. Como a taxa de evasão usa matrículas elegíveis no denominador, a duplicação altera contagem e perfil do grupo. A secretaria precisa corrigir o grão antes de agregar. O parecer acadêmico informa que a transferência não equivale a uma segunda exposição ao risco: ela muda o vínculo administrativo dentro do mesmo semestre. A trilha de reconciliação deve conservar os dois registros de origem, selecionar o estado válido e provar unicidade na tabela analítica. O relatório será auditado por curso e turno, portanto uma exclusão aleatória poderia alterar justamente os estratos menores. A regra deverá usar chave composta, vigência e estado acadêmico, registrar conflitos e interromper a publicação quando mais de uma linha permanecer elegível para a mesma chave. Só depois desse teste o numerador de evasão e o denominador de estudantes expostos poderão ser agregados sob a mesma definição.

Essa definição institucional vigente também deve constar da linhagem publicada.

\FonteDidatica

\Comando Selecione a ação que restaura uma linha por estudante-semestre e protege numerador e denominador.

\begin{enumerate}[label=\Alph*)]
  \item somar todas as linhas porque representam eventos reais
  \item definir chave estudante-semestre, reconciliar transferências e testar unicidade antes da taxa
  \item usar a média das taxas publicadas
  \item remover aleatoriamente uma duplicata
  \item agregar por curso antes de investigar
\end{enumerate}

\end{document}
